\documentclass[11pt]{article}
\usepackage[margin=1in]{geometry}
\usepackage{amsmath,amssymb,mathtools,url}
\setlength{\parindent}{0pt}
\setlength{\parskip}{5pt}
\setlength{\emergencystretch}{2em}
\DeclareMathOperator{\Tr}{Tr}
\begin{document}
\paragraph{Strand X: bounded native Gaussian preparation result
[N/S; V:X,W,D].}
The declared zero-coupling control gives two accepted intermediate
$5^3$ preparation states, with small floating-point clock and ruler
residuals and positive fields inside the stated point guards.
It reaches none of the prescribed loads and performs no time step.
The $7^3$ parent passes its flat controls and retains only its flat
state when the whole-batch wall cap is detected.
The numerical content [N] is this finite-parent preparation and
stored-array readout, not an enclosed root or a trajectory.
The exact residual identities and the interpretation below have
supporting [S] scope.  No new [X] or [R] claim is assigned.
The single independent result review
\path{solve/review_X_gaussian_result.tex} passes the actual source,
stored-output interpretation and two-sentence delivery at this
preparation-only scope, with no additional scientific repair.

\paragraph{1. Fixed law, state and acceptance meaning.}
The implemented control is XN11--XN16 of
\path{solve/strand_X_nonlinear_response.tex}, with
\[
 a=1,\qquad a^2m_s^2=(1/12,1/6,1/6,1/4),\qquad
 \eta=(+,-,-,+),\qquad g_M=g_B=0 .
\]
The prescribed $\Lambda=m_1$ and $\ell=a$ do not enter the
zero-coupling seed.  The single coefficient match is retained;
no radius, mass, lapse or Newton coefficient is fitted to the output.
Each Cartesian parent uses the open native sine stencil,
all interior--collar links and all native modes.  The outermost
layer has $N_B=\psi_B=1$.  There are respectively 27 and 125
interior rulers, 250 and 686 native two-spin modes, and 500 and
1372 carrier modes per fermion packet.
All four actual fermionic seas are retained.  The invariant empty
sector of the untruncated bosons retains the original vacuum scalar.
There is no spherical averaging or interacting-light approximation
inside this selected Gaussian control.

The light covariance is the specified
$C_1=P_{1,-}^{0}+\lambda P_g$, with $P_g$ the fixed normalized
rank-two positive-band preparation, and the other covariances are
their flat seas.  The complete native state is held fixed during
each initial field solve; its mismatch with the sea at the solved
ruler is included in the sources.  The prescribed loads are
$1/256,1/16,1$.  A fractional intermediate $\lambda$ is a mixed
Gaussian preparation label with added mean charge $2\lambda$;
it is neither physical time nor a replacement target load.

The complete rectangular Jacobian has rows on all $\Lambda$ and
ruler columns only on $I$:
\begin{equation}
 C_x=w_x+\Tr(\rho e_x),\qquad
 J_{xy}=\partial_{\psi_y}C_x\big|_\rho,\qquad
 J_{II}^{T}N_I=-J_{BI}^{T}{\bf1}_B .
 \tag{XR1}
\end{equation}
The code retains these collar rows and the actual polar derivatives.
Its clock solve uses $10^{-12}$, stricter than the declared
$10^{-10}$ clock acceptance tolerance; the ruler tolerance remains
$10^{-10}$.  The prescribed trial guards are
\[
 1/8<N_I<8,\qquad 1/2<\psi_I<4,\qquad
 \sigma_{\min}(J_{II})/\sigma_{\min}(J_{0,II})>10^{-4}.
\]
Every Newton/backtracking evaluation consumes the same 12-trial
field-solve allowance and the whole-batch field counter.
A completed field evaluation need not be an accepted clock root.
In particular the lapse guards are applied at trial fields before
the clock has reached its tolerance.
All numerical residuals, singular values and guard tests here are
ordinary binary64 diagnostics, without directed root, inverse or
trajectory error bounds.

\paragraph{2. What is accepted at each parent.}
Both flat controls pass their declared numerical tolerances:
the individual sea clock and derivative rows vanish to roundoff,
the full vacuum column sum vanishes to roundoff, and the
constant-interior, central-site and alternating-interior quadratic
forms agree with the independent spectral commutator formula.
The column-sum infinity norms are
$2.153565991234155\,10^{-14}$ for $5^3$ and
$2.1485408717864798\,10^{-14}$ for $7^3$.
The recorded quadratic-form differences are at most
$7.147060721024445\,10^{-15}$ and
$6.619704784327496\,10^{-15}$, respectively.
The code's field named \texttt{relative\_difference} divides by
$\max(1,|\hbox{left}|,|\hbox{right}|)$; these are differences with
that reporting floor, not relative errors divided by a small
quadratic form.  Three tested profiles do not constitute an
all-field numerical coercivity theorem.

The $5^3$ accepted preparations are
\[
\begin{array}{c|c|c|c}
 \lambda&\|C_I\|_\infty&\|J^TN\|_\infty&
                         \sigma_{\min}(J_{II})\\ \hline
 1/16384&4.163447049550543\,10^{-14}&
          1.0551569912460471\,10^{-17}&0.005727854876008834\\
 5/65536&4.501707274756384\,10^{-14}&
          1.3553400339841182\,10^{-17}&0.004365295249808526
\end{array}
 \tag{XR2}
\]
At $1/16384$ the interior lapse lies between
$1.2854624365981677$ and $2.098965707565491$.
Its field record survives, but its full state arrays were
overwritten by the retained preparation checkpoint.
No later array readout is assigned to that earlier state.

The actual retained $5^3$ checkpoint has
\begin{equation}
\begin{aligned}
 \lambda&=5/65536,&
 \psi_I&\in[1.0003327879495905,\;1.0023484106759575],\\
 N_I&\in[1.5170668299235484,\;3.1315020428088975],&
 \frac{\sigma_{\min}(J_{II})}{\sigma_{\min}(J_{0,II})}
   &=0.4946556144874843 .
\end{aligned}
\tag{XR3}
\end{equation}
Its marked trace is $0.000152587890625=2\lambda$.
The stored covariance guard reports maximum spectral
$[0,1]$ defect $5.10702591327572\,10^{-15}$ and maximum
Hermiticity Frobenius defect $1.8842851490861374\,10^{-22}$.
The retained-array reader did not recompute covariance spectra.
The stored zero spectrum-conservation defect at preparation is
not a test of dynamical conservation.

The $7^3$ retained checkpoint is $\lambda=0$, with $\psi_I=1$,
$N_I\in[0.999999999998313,1.0000000000005618]$,
clock residual $1.2603146094162877\,10^{-14}$, ruler residual
$1.3613513528271115\,10^{-17}$ and
$\sigma_{\min}(J_{0,II})=0.005773665332074604$.
There is no common positive accepted load between the two parents,
so the data supply no finite-parent convergence comparison.

\paragraph{3. Exact accounting identities and their numerical defects.}
Let
\[
 E_{\rm col}=\sum_{b\in B}C_b,\qquad
 E_{\rm pos}=\sum_{x\in\Lambda}\frac{\Tr(\rho e_x)}{\psi_x},
 \qquad {\cal H}=\sum_xN_xC_x .
\]
The finite polar trace gives $\sum_xw_x/\psi_x=0$ exactly.
At an arbitrary field, without assuming $C_I=0$, the two
distinct identities are
\begin{equation}
 E_{\rm col}-E_{\rm pos}=-\sum_{i\in I}C_i/\psi_i,\qquad
 {\cal H}-E_{\rm col}=\sum_{i\in I}N_iC_i .
 \tag{XR4}
\end{equation}
At an exact root these reduce to the positive charge identity
of XN5.  At the retained $5^3$ checkpoint the stored-array
readout gives
\[
\begin{array}{c|r}
 E_{\rm col}&0.00026815188606647806\\
 E_{\rm pos}&0.0002681518861636053\\
 E_{\rm col}-E_{\rm pos}&-9.712722163005005\,10^{-14}\\
 -\sum_I C_i/\psi_i&+1.3435433109824023\,10^{-13}\\
 {\cal H}-E_{\rm col}&-3.4801827218655657\,10^{-13}\\
 \sum_I N_iC_i&-3.4801829438392725\,10^{-13}
\end{array}
\]
Thus the first identity's signed numerical defect is
$-2.314815527282903\,10^{-13}$, and the second's is
$2.219737068474066\,10^{-20}$.  The first is not set to zero
or hidden by the small root residual: in exact arithmetic its
remaining term is the finite polar trace sum.  These are
single-state accounting diagnostics, not charge conservation
measurements over time.
At the flat $7^3$ checkpoint the reported
$E_{\rm col}=4.950707796706056\,10^{-13}$ and
$E_{\rm pos}=7.971047946312673\,10^{-13}$ are numerical
residuals of the nominal zero vacuum charge, not positive
physical energies.

The same retained rectangular matrix supplies the exact algebra
\begin{equation}
 r_1=J_{\Lambda I}^{T}{\bf1},\qquad
 J_{II}^{T}(N_I-{\bf1})=-r_1,\qquad
 \|N_I-{\bf1}\|_\infty
       \leq\frac{\|r_1\|_2}{\sigma_{\min}(J_{II})}
       =:\beta .
 \tag{XR5}
\end{equation}
At $5^3$ the separately extracted transpose-identity residual
is $3.604250708511724\,10^{-17}$ and
$\beta=6.7169792935201$, whereas the actual
$\|N_I-{\bf1}\|_\infty=2.1315020428088975$.
A certified $\beta<7/8$ would suffice for both declared lapse
guards.  Its failure here is inconclusive; it is not a lapse
or branch failure, and the actual point passes both guards.
Neither the singular value nor this quotient is a directed
neighborhood certificate.

\paragraph{4. Separate the trial limit, step floor and wall cap.}
The $5^3$ terminal continuation attempt starts from $5/65536$
and proposes $3/32768$ with increment $1/65536$.
It exhausts the 12-trial field solve.  The next halved
increment is $1/131072$, strictly below the declared
$2^{-16}$ floor, and preparation stops.
The attempted increment itself equals the floor.
The terminal completed trial has clock residual
$1.1597454869122868\,10^{-7}$, so it is not an accepted root,
even though its interior lapse maximum
$7.965514209890777$ is below eight.
Other rejected trials at that same candidate have maximum
lapse above eight; they likewise have unsolved clock rows.
These records do not locate a physical lapse zero, a fold,
or the boundary of the exact positive branch.

The $7^3$ parent attempts only the first prescribed load
$1/256$ from its flat state.
Trial fields encounter the lower lapse guard; one encounters
both lower and upper guards.  Their clock residuals remain
nonzero.  Field request 125, at backtracking index four,
ends at the wall cap before completing the source/Jacobian
and lapse readout.  No missing lapse or singular value is
inferred for this request.  There is no completed $7^3$
halving sequence or step-floor outcome.
The TXT's terminal corrector reason \texttt{pending} is an
incomplete row; the final JSON top-level failure identifies
the actual \texttt{Cap} phase.

The final status is \texttt{unresolved}.
The authoritative JSON/TXT elapsed value is
$21628.5037657097$ seconds, exceeding the declared
$21600$ seconds by $28.5037657096982$ seconds.
The subsequent log samples $21628.560513593256$ seconds.
The two finalization samples are retained literally.
The cooperative cap detection does not establish strict
six-hour wall compliance.
The cumulative counters are 125 requested field/source
evaluations, comprising 46 completed and 79 failed or
incomplete evaluations, 1182 scalar integrand evaluations,
and zero accepted time steps.
The actual backend is CPU at nice 10 with seven own threads;
reported peak host RSS is 2236080128 bytes and device-pool
allocation is zero.  No GPU execution is credited.
The measured flat-field durations are
$13.024783998727798$ seconds for $5^3$ and
$2916.5891262441874$ seconds for $7^3$; these are this
implementation's recorded costs, not a complexity theorem.

\paragraph{5. Source/readout boundary and missing information.}
The original interrupted run retains its JSON/TXT/log and
35 charged field requests but no NPZ checkpoint.
The reporting repair reconstructs the exact declared flat
input, retains the prior counters and repair downtime, and
saves the actual last-valid states.  It changes no physical
input, branch guard, method or budget.
The current model source contains eight prospective
\texttt{Parent.assemble} lines computing the $r_1$,
transpose-identity and $\beta$ report.  Those lines were
added after the continued process launched and did not
execute in either model run.
The values in (XR5) are the separate existing-data reader's
stored-array readout, using the matching stored singular value;
they do not retrospectively become model diagnostics.

The single data-only reader uses no new field/source call,
eigensolve or evolution.  It records one thread, nice 10,
$0.05470031499862671$ seconds and 61337600 bytes peak RSS.
Rejected trial matrices and overwritten earlier accepted
arrays are unavailable.  The optional load-affine
flat-start comparison was not performed because its
required pairs of full matrices were not retained.
No inference in this result requires reconstructing them.

\paragraph{6. Narrowed conclusion and constructive handoff.}
The exact finite-parent local branch theorem of XN1--XN10
is consistent with the two numerical preparations.
Its implicit-function neighborhood was not quantitatively
identified with any prescribed load.  A pointwise positive
Jacobian singular value and small residual do not alone
give a root enclosure or continuation radius.
The present data therefore distinguish a realized
small-load preparation from an unresolved target
preparation; they do not refute the exact field equations
or establish failure of the physical collapse law.

The next solver ingredient is the branch-local
residual/inverse/Jacobian-variation control already
specified in XN9, applied to the actual fixed-native
sources with their collar derivatives.
For example, if $M$ is a bounded approximate inverse at
a declared candidate $\psi_*$, an enclosure of
\[
 \sup_{\psi\in{\cal B}_r}
      \|I-MJ_{II}(\psi)\|\leq q<1,\qquad
 \|M C_I(\psi_*)\|+qr\leq r
\]
together with lapse and ruler bounds on the same ball
would distinguish a certified root neighborhood from a
failed numerical corrector.  The available stored point
does not supply that ball bound.  The large measured
$7^3$ full-derivative cost also has to be addressed before
an exact-source trajectory is a practical test: tangent
and transpose operations through the same polar Sylvester
derivatives are possible implementation directions, but
no replacement solver or new run is adopted here.
No branch guard, load target or continuation floor is
relaxed to manufacture an accepted result.

An actual collapse handoff still requires a prescribed
load prepared on the declared branch, both state and
constraint evolution with endpoint field solves, and the
same-history charge, material-radius, native-light-speed,
lock-defect and finite-frequency source monitors of X/XS/XN.
The frozen-field susceptibility is only a component of
the coupled response, not its clock/lapse-eliminated value.
The native light-algebra entropy and cross-surface
correlations of XP use that same trajectory and declared
cut.  None is supplied by a preparation checkpoint.
There is no time-step comparison, causal sea response
measurement, charge-fixed shrinkage, spherical/open
transfer, production spectrum, Page curve or thermal
law in this output.  Nonzero dressed couplings still
require the complete light sources and actual state.
This handoff identifies the unresolved ingredient;
it authorizes no additional computation.

\paragraph{Paper delivery.}
At the declared zero-coupling native control, the $5^3$ parent
admits two numerically accepted intermediate preparations,
the retained $\lambda=5/65536$ state having clock residual
$4.50\,10^{-14}$ and positive interior lapse
$1.517<N_I<3.132$, without a root enclosure.
Neither parent reaches a prescribed load or evolves in time,
so the step-floor and wall-cap outcomes leave collapse,
causal response, Page correlations and thermal production
uncomputed.

\paragraph{SOURCES.}
The governing analytical system and fixed control are
XN1--XN16 in \path{solve/strand_X_nonlinear_response.tex}
and its single review
\path{solve/review_X_nonlinear_response.tex}.
The causal and readout dependencies are
\path{solve/strand_X_causal_sea.tex},
\path{solve/strand_X_page_correlations.tex}, and
\path{solve/strand_W.tex}, at their stated scopes.
The actual model source is
\path{solve/code/x_nonlinear_gaussian.py}, with the
unexecuted eight-line prospective metric delta distinguished
above; the actual retained-array reader is
\path{solve/code/x_nonlinear_existing_data.py}.
The unchanged original and continued JSON/TXT/log artifacts
are in \path{solve/results/x_nonlinear_gaussian/}, under
the stems \texttt{x\_nonlinear\_gaussian} and
\texttt{x\_nonlinear\_gaussian\_continued}.
The complete literal assessment is
\path{solve/results/x_nonlinear_gaussian/x_nonlinear_gaussian_existing_data.txt}:
final status and sampling at lines 1--60,
$5^3$ controls/attempts at 61--456,
retained $5^3$ arrays at 457--499,
$7^3$ controls/attempt at 762--840,
retained $7^3$ arrays at 841--883,
trial records and missing-data/resource statements thereafter.
The two actual checkpoint paths are
\path{solve/results/x_nonlinear_gaussian/L5_last_valid_preparation.npz}
and
\path{solve/results/x_nonlinear_gaussian/L7_last_valid_preparation.npz},
also retained under
\path{<internal>/}
on the scientific host; they are not Git payloads.
The independent result assessment is
\path{solve/review_X_gaussian_result.tex}.
The [N] claims are bounded floating-point measurements;
the residual and branch implications are supporting [S]
uses of the retained analysis, not new independent-method
or physical results.
\end{document}
